\chapter{Utility Scripts}
\label{ch:scripts}

The \file{scripts/} directory holds stand-alone tools built on the
\code{jnuscout} package: uncertainty calibration, active-learning data
preparation, database hygiene, and structure checks.  The scripts that use
the package embed a small \code{sys.path} bootstrap for the source layout, so
they can be run directly from a source checkout:

\begin{transcript}
$ python scripts/<script>.py --help
\end{transcript}

The \code{--help} output of a script is its authoritative option list; this
chapter gives the purpose, the dependencies, a canonical invocation and the
output of every tool.  Scripts that talk to a SCINE database additionally
need \code{pymongo} and the \code{[scine]} extra; scripts that label or
validate with ORCA need an ORCA executable (\code{ORCA_BIN}); the RDKit
structure check needs a second interpreter that has RDKit installed.
Most tools write their output under \file{results/} or \file{data/}, and
inputs are not modified; the exceptions are the two reporting tools
(\script{analyze_committee} and \script{db_snapshot}), which print to the
terminal, and the database tools, which write to the database they are
pointed at (\code{setup_h2co_nnp_db --wipe} deletes it first).

\begin{table}[htbp]
\centering
\caption{Overview of the utility scripts.  ``Needs'' lists the dependencies
beyond the Python standard library.}
\label{tab:scripts-overview}
\begin{tabularx}{\textwidth}{@{}lXl@{}}
\toprule
Script & Purpose & Needs \\
\midrule
\script{calibrate_uq} & calibrate the fallback threshold and audit it & numpy, ORCA \\
\script{build_al_trainset} & sample reaction paths, label frames with ORCA L2 & pymongo, SCINE database, ORCA \\
\script{resplit_trainset} & composition-stratified, reaction-grouped re-split & --- \\
\script{validate_finetune} & held-out barrier validation of a fine-tuned model & pymongo, MACE, ORCA \\
\script{analyze_committee} & member and combination statistics on a barrier file & numpy \\
\script{db_snapshot} & progress snapshot of the search databases & SCINE database \\
\script{check_structure_identity} & independent RDKit cross-check of database structures & SCINE database, DFTB3, RDKit \\
\script{eval_barrier_spline} & spline-endpoint barrier evaluation with ORCA caching & pymongo, SCINE database, ORCA \\
\script{setup_h2co_nnp_db} & build and run the H$_2$CO test database with the NNP family & SCINE database, Puffin \\
\script{verify_bohr_units} & known-answer check of the Bohr-to-Angstrom conversion & pymongo, ORCA \\
\script{_rdkit_perceive} & perception helper for the identity check (not run directly) & RDKit \\
\bottomrule
\end{tabularx}
\end{table}

\section{\texttt{calibrate\_uq}}
\scriptlabel{calibrate_uq}

Builds a reference set of small CHNO molecules plus random displacements,
computes the committee prediction and disagreement and the ORCA L2 reference
for every configuration, and calibrates the fallback threshold with
\code{calibrate_threshold}: the low quantile of the disagreement among the
configurations whose true error exceeds the danger threshold.  It also
reports the rank correlation, the expected calibration error and a coverage
scan, and prints the acceptance targets for the user to judge (rank
correlation above 0.6, coverage above 80\%, fallback rate below 20\%).

\begin{description}
  \item[Dependencies] \code{numpy}; the NNP committee (the \code{[nnp]}
        extra, including torch and the model weights); ORCA for the
        reference calculations, unless \code{--no-orca} is given.
  \item[Usage] \code{python scripts/calibrate_uq.py --n 60 --epsilon 0.10 --out results/uq_calib.json}
  \item[Options] \code{--n} displaced configurations per molecule (5);
        \code{--sigma} displacement width in \AA{} (0.05, near-equilibrium
        thermal noise); \code{--epsilon} danger threshold as per-atom force
        error (0.10); \code{--theta-quantile} (0.05); \code{--device}
        (cuda); \code{--limit} to restrict a debug run; \\
        \code{--calibrate-baselines} to align the per-member energy
        baselines first; \code{--no-orca} to measure the committee alone,
        which yields a disagreement distribution but no calibrated
        threshold.
  \item[Output] \file{results/uq_calib.json} with the summary (threshold,
        rank correlation, expected calibration error, coverage statistics),
        the per-configuration rows and the failures; the summary is also
        printed.
\end{description}

\section{\texttt{build\_al\_trainset}}
\scriptlabel{build_al_trainset}

Pulls elementary steps from a SCINE database, samples configurations along
each fitted path with \code{spline.evaluate(t)}, labels them with ORCA L2
(r2SCAN-D4/def2-SVP) in parallel worker processes, and writes the frames for
MACE fine-tuning in extended-xyz format, in eV and \AA.  Sampling is
transition-state weighted by default.  \code{--holdout-es} reserves a number
of elementary steps as the held-out set, split by step so that no path point
leaks between the two sides.

\begin{description}
  \item[Dependencies] \code{numpy}, \code{pymongo}, the SCINE database, and
        ORCA at level L2.
  \item[Usage] \code{python scripts/build_al_trainset.py --db diels_alder --n-points 15 --workers 8 --out data/al_trainset.xyz}
  \item[Options] \code{--db} (a database name, or several separated by
        commas; default \code{diels_alder}); \code{--n-points} per
        elementary step (15); \code{--max-es} upper bound on the steps
        (40); \code{--workers} parallel labellers (8); \code{--out}
        (\file{data/al_trainset.xyz}); \code{--holdout-es} number of steps
        to reserve (0); \code{--holdout} to generate the held-out set
        itself; \code{--uniform} for uniform instead of transition-state
        weighted sampling.
  \item[Output] one extxyz file; each frame carries the energy, the forces
        and the bookkeeping fields (\code{es_id}, \code{t}).  Frames whose
        reference calculation fails are skipped and counted.
\end{description}

\section{\texttt{resplit\_trainset}}
\scriptlabel{resplit_trainset}

Merges one or two existing frame files and re-splits them stratified by
composition.  Within each composition the elementary steps are shuffled and
divided 8:2, and all frames of one step go to the same side, so that
consecutive geometries of a path never straddle the split.  It exists because
an order-based split produced a training set entirely of C$_4$H$_6$ and a
validation set entirely of C$_6$H$_{10}$.

\begin{description}
  \item[Dependencies] the Python standard library only.
  \item[Usage] \code{python scripts/resplit_trainset.py --a data/al_train_da.xyz --b data/al_holdout_da.xyz --out-train data/al_train2.xyz --out-valid data/al_valid2.xyz}
  \item[Options] \code{--a} the primary frame file; \code{--b} an optional
        second file to merge in (duplicate geometries are dropped);
        \code{--out-train} and \code{--out-valid} the output files;
        \code{--valid-frac} (0.2); \code{--seed} (20260911).
  \item[Output] two extxyz files, plus the per-composition frame and step
        counts for both sides, so the composition balance can be checked
        before training.
\end{description}

\section{\texttt{validate\_finetune}}
\scriptlabel{validate_finetune}

Evaluates a fine-tuned model against the foundation model and the ORCA L2
reference on held-out elementary steps, and reports the barrier mean absolute
error together with its bias decomposition.  The held-out steps are read from
the \code{es_id} fields of the validation file, which guarantees that no path
point of the training set is reused.

\begin{description}
  \item[Dependencies] \code{numpy}, \code{pymongo} (structure lookup in the
        database), MACE (both the fine-tuned and the foundation model), and
        ORCA.
  \item[Usage] \code{python scripts/validate_finetune.py --valid data/al_valid2.xyz --finetuned models/mace_al_da.model --baseline <foundation.model> --db diels_alder --out results/finetune_eval.json}
  \item[Options] \code{--valid} the validation frames; \code{--finetuned}
        the fine-tuned checkpoint; \code{--baseline} the foundation
        checkpoint, which is required; \code{--db}; \code{--device}
        (cuda); \code{--out} (\file{results/finetune_eval.json}).
  \item[Output] the JSON file (summary and rows) plus the printed summary:
        barrier MAE and bias per model, and the per-reaction rows the
        decomposition is built from.
\end{description}

\section{\texttt{analyze\_committee}}
\scriptlabel{analyze_committee}

Compares committee member combinations on a barrier file produced by the
evaluation scripts.  For every single member and every equal-weight
combination it reports the sample size, MAE, bias, RMSE, rank correlation and
slope, so that the committee design can be judged on corrected-unit
geometries instead of intuition.  Rows whose reference lies below a
threshold are excluded, because barrierless or unphysical entries would
distort the statistics.

\begin{description}
  \item[Dependencies] \code{numpy}.
  \item[Usage] \code{python scripts/analyze_committee.py --in results/barrier_committee.json}
  \item[Options] \code{--in} the input barrier file (default
        \file{results/barrier_committee.json}); \code{--neg-thresh} drop
        rows with an ORCA reference below this value (default $-1.0$
        kcal/mol).
  \item[Output] printed tables of the single-member and combination
        statistics; no file is written.
\end{description}

\section{\texttt{db\_snapshot}}
\scriptlabel{db_snapshot}

Refreshes the summary numbers of the search databases.  For each database it
prints the structure count, the elementary-step count, the size histogram of
the structures and the reaction-type histogram of the elementary steps; this
is the snapshot used for progress reporting.

\begin{description}
  \item[Dependencies] the SCINE database (a running MongoDB).
  \item[Usage] \code{python scripts/db_snapshot.py --dbs diels_alder_nnp h2co_nnp}
  \item[Options] \code{--dbs} one or more database names (defaults
        \code{diels_alder_nnp} and \code{h2co_nnp}).
  \item[Output] stdout.  Structures whose size cannot be read and steps
        whose type cannot be classified are reported as counts, not silently
        dropped.
\end{description}

\section{\texttt{check\_structure\_identity}}
\scriptlabel{check_structure_identity}

Independent-source cross-check of database structures.  For each sampled
structure it writes an XYZ file, lets RDKit (in a separate interpreter)
perceive connectivity and bond orders from the coordinates, regenerates a 3D
reference from the resulting SMILES, relaxes that reference on the same DFTB3
surface and compares the energies; a clearly positive difference means the
library entry is not the lowest configuration of that molecule.  Statistics
must be grouped by label: the default filter considers only structures that
claim to be relaxed minima, because transition states should not sit at
minima by definition, so a large difference there is expected rather than a
defect.

\begin{description}
  \item[Dependencies] the SCINE database, DFTB3 (Sparrow) for the
        relaxation, and RDKit in a separate interpreter.  Set
        \code{JNUSCOUT_RDKIT_PY} to an interpreter that has RDKit
        installed; the helper defaults to \script{_rdkit_perceive} next to
        the script and can be overridden with \code{JNUSCOUT_RDKIT_HELPER}.
  \item[Usage] \code{python scripts/check_structure_identity.py --db diels_alder --label minima --sample 60}
  \item[Options] \code{--db}; \code{--id} to check a single entry;
        \code{--sample} (40); \code{--flag} the flagging threshold in
        kcal/mol (5.0); \code{--label} one of \code{minima}, \code{guess},
        \code{saddle}, \code{all} (default \code{minima}); \code{--size}
        atom-count filter; \code{--group} for grouped statistics;
        \code{--out} (\file{results/identity_check.json}).
  \item[Output] the JSON file (parameters, exclusion reasons, rows) and the
        printed summary.  Trusted counts and exclusion reasons are reported
        together, because a percentage without them is misleading.
\end{description}

\section{\texttt{eval\_barrier\_spline}}
\scriptlabel{eval_barrier_spline}

Barrier evaluation by the spline-endpoint method: the reactant is
\code{spline.evaluate(0)} and the transition state is
\code{spline.evaluate(ts_position)}, so the reference and every model see
exactly the same two geometries.  This is the correct protocol for
bimolecular elementary steps, where a path cannot be assembled by
concatenating the reactant fragments.  Several models can be evaluated in one
run, and ORCA results are cached across runs.

\begin{description}
  \item[Dependencies] \code{numpy}, \code{pymongo}, the SCINE database,
        ORCA, and the MACE models to compare.
  \item[Usage] \code{python scripts/eval_barrier_spline.py --es-file data/holdout_bi.json --model v3f=models/mace_al_v3f.model --out results/barrier_bi.json}
  \item[Options] \code{--es-file} (required: a specification with the
        database name and the elementary-step ids); \code{--db} to override
        the database; \code{--model name=path} (repeatable);
        \code{--device} (cuda); \code{--workers} (12); \code{--orca-cache}
        (\file{results/barrier_spline_orca_cache.json}); \code{--out}
        (\file{results/barrier_spline.json}).
  \item[Output] the barrier JSON (summary and rows) and the ORCA cache; the
        cache-hits and the number of ORCA evaluations are printed, so a
        resumed run is visibly cheaper.
\end{description}

\section{\texttt{setup\_h2co\_nnp\_db}}
\scriptlabel{setup_h2co_nnp_db}

Builds a Chemoton database whose model carries \code{method_family='nnp'},
following the H$_2$CO assembly (the same gears and the same AFIR job as the
reference setup), and optionally submits one exploration cycle.  The
coordinates are defined in Angstrom and converted to Bohr before insertion,
because the database stores structures in atomic units.  Puffin, with the NNP
patch enabled, must already be running (Chapter~\ref{ch:scine-integration}).

\begin{description}
  \item[Dependencies] \code{numpy}, the SCINE database and utilities, a
        running MongoDB, and a patched Puffin daemon.
  \item[Usage] \code{python scripts/setup_h2co_nnp_db.py --run}
  \item[Options] \code{--run} to run one cycle after assembly; \code{--wipe}
        to wipe the database first; \code{--only-run} to skip the insertion
        and run the gears for \code{--loops} rounds; \code{--loops} (1);
        \code{--db} (\code{h2co_nnp}).
  \item[Output] the database contents and, with \code{--run}, one submitted
        exploration cycle.  Housekeeping needs several gear rounds before
        compounds appear, which is what \code{--loops} is for.
\end{description}

\section{\texttt{verify\_bohr\_units}}
\scriptlabel{verify_bohr_units}

Checks that geometries read from the search database are converted from Bohr
to Angstrom correctly.  The script evaluates the same elementary step under
both conventions and quantifies the effect on the ORCA barrier, and it
performs a known-answer self-check: the nearest-neighbour distance of an
n=6 reactant must come out as a normal C--H bond length, not a broken-bond
distance.  This is the reproduction tool for the unit defect that once
stretched every geometry by a factor of 1.8897.

\begin{description}
  \item[Dependencies] \code{numpy}, \code{pymongo}, and ORCA.
  \item[Usage] \code{python scripts/verify_bohr_units.py --es <es-id> --out results/bohr_unit_check.json}
  \item[Options] \code{--db} (\code{diels_alder}); \code{--es} (required:
        the 12-character elementary-step id prefix); \code{--model} an
        optional extra model to evaluate; \code{--out}
        (\file{results/bohr_unit_check.json}).
  \item[Output] the JSON file with both barrier values and the distance
        check; the two barriers should differ by much more than chemical
        accuracy when the conversion is wrong, which is the whole point.
\end{description}

\section{\texttt{\_rdkit\_perceive}}
\scriptlabel{_rdkit_perceive}

Internal helper for \script{check_structure_identity}; it is not meant to be
run directly.  It perceives connectivity and bond orders from 3D coordinates
with RDKit's \code{rdDetermineBonds}, writes the canonical SMILES, and
regenerates a 3D reference structure from it.  Charged fragments are
disabled, so open-shell structures are represented with radical electrons
instead of fragment ions, which is the correct choice for the neutral systems
the check targets.

\begin{description}
  \item[Dependencies] RDKit, in the interpreter selected by the caller.
  \item[Usage] \code{<rdkit-python> _rdkit_perceive.py <in.xyz> <out_prefix>}
  \item[Output] \code{<out_prefix>.smi} (the SMILES; the first stdout line
        carries the same text), \code{<out_prefix>.xyz} (the regenerated
        structure), and a diagnostics line on stderr reporting formal
        charges, radical electrons and the fragment count, so the caller can
        judge whether the perception is trustworthy.
\end{description}
